% !TEX root = ../main.tex

% 中英标题：\chapter{中文标题}[英文标题]
\chapter{课题主要研究内容及进度}


\section{课题主要研究内容}

蛋白质是生命系统的重要组成部分，深入理解蛋白质功能是揭示生物系统运作机理、解析分子尺度的生物行为以及探索蛋白质多样性和进化关系的关键环节\cite{eisenberg2000protein}。此外，精确预测蛋白质功能在阐明基因功能、研究基因调控网络、指导未表征蛋白质的功能注释以及促进药物发现和开发等领域均具有重大意义。准确的功能预测能大幅提高药物开发的针对性和疗效，从而进一步推动医学进步，改善人类健康与生活质量。

然而，全球最权威的蛋白质数据库UniProt\cite{uniprot2021uniprot}中，经过实验标注的蛋白质比例不到1\%，即使在人类蛋白质组中，超过三分之二的蛋白质功能也仍然未完全标注或尚未标注\cite{boutet2016uniprotkb}。因此，自动化的蛋白质功能预测方法近年来逐渐受到研究者的高度关注。相比传统的实验手段，计算方法更高效、经济且具有强大的扩展性和广泛的适用性，逐渐成为当前研究的热点。

随着高通量生物数据的爆炸式增长，基于单一模态（如序列或结构信息）的传统功能预测方法在实际应用中已逐渐暴露出明显的不足。这类方法在处理标签稀疏、功能类别分布不均衡、跨物种数据差异巨大等复杂情境时往往表现欠佳。近年来，多模态功能预测方法开始兴起，研究者探索融合蛋白质序列、结构以及功能标签语义信息，试图通过多模态协同建模获得更通用和更具语义感知能力的模型\cite{wang2023large}。然而，现有多模态方法仍存在以下突出问题：

(1) 模态融合方法粗浅，缺乏统一高效的表示空间: 大多研究仅以简单拼接或并行网络方式整合不同模态的特征，无法充分发挥模态间的深层交互和协同作用，限制了模型在任务迁移和零样本预测中的表现。

(2) 过度依赖监督训练，泛化能力不足:现有方法多将功能预测视为标准多标签分类任务，严重依赖于已标注数据的标签分布。尤其在面对低资源或长尾标签情景时，模型泛化性能明显下降。此外，多数方法未有效利用功能标签的语义信息，缺乏对功能语义结构的深层理解，进一步限制模型的泛化能力。

针对上述问题，本课题提出了一种新型的多模态蛋白质功能预测框架，通过融合蛋白质的序列特征、结构信息以及功能标签语义表示，并结合有监督学习与对比学习的多任务学习方法，构建一个统一的蛋白质功能预测模型。本课题的具体研究内容包括以下方面：

(1) \textbf{数据获取与预处理：} 本研究所需数据包括蛋白质的序列数据、结构数据和功能标签，需要对这些不同来源的数据进行精准对齐和筛选，选取具有专家手工注释和实验验证的高可信功能标签。

(2) \textbf{蛋白质多模态信息的预训练特征建模：} 序列方面，本研究采用高效预训练语言模型ESM-2\cite{doi:10.1073/pnas.2016239118}对氨基酸序列进行编码，提取丰富的序列语义信息。结构方面，利用Cα原子坐标构建蛋白质接触图，设计了序列边、K近邻边和半径边三种边类型，以全面捕捉蛋白质结构的空间拓扑特征与物理邻接关系。基于此构建异构图，充分利用异构图神经网络（HGNN\cite{zhang2019heterogeneous}）的强大表达能力，有效融合蛋白质序列与结构信息，生成统一且高质量的蛋白质多模态表示。

(3) \textbf{多任务学习框架设计：} 为有效提升模型的预测精度与语义对齐能力，本研究设计了一种创新的多任务学习框架，首先在自监督预训练阶段进行特征重建与掩码恢复任务，进一步强化模态表示的鲁棒性和泛化性。在有监督训练阶段，通过先进的文本嵌入模型Qwen3-Embedding\cite{qwen3embedding}为功能标签生成语义嵌入，将其与蛋白质多模态表示通过交叉注意力与拼接方式融合，实现深入且细致的语义交互与融合。主任务采用MLP分类器\cite{taud2017multilayer}，并以Focal Loss\cite{lin2017focal}优化基因本体论（GO）功能预测，增强对长尾和低频标签的识别能力。辅助任务则基于SupCon Loss\cite{khosla2020supervised}进行蛋白质与功能标签之间的有监督对比学习，有效提升蛋白质表示与功能语义空间的匹配精度和泛化能力，从而为零样本预测任务奠定坚实基础。

通过以上研究，本课题期望显著提升蛋白质功能预测的准确性、泛化能力及对新功能类别的预测能力，为蛋白质功能预测领域提供更有效的计算工具和理论支持。




\section{进度介绍}

截至目前，本课题已顺利完成了数据准备和模型搭建的关键阶段，奠定了完整研究工作的基础。在数据获取和预处理阶段，已成功从数据库中获取了大量蛋白质序列、结构信息和功能标签数据;在多模态特征建模方面，完成了序列与结构信息的有效编码与融合;在功能语义建模阶段，课题已引入大语言文本嵌入模型Qwen3-Embedding为GO功能描述生成高质量的文本语义嵌入，使功能标签具备清晰、丰富的语义表示，从而有力地支持蛋白质-功能语义匹配任务。课题搭建了一个多任务联合训练框架，通过交叉注意力与拼接方式，实现了蛋白质多模态表示与功能语义表示的深度融合。

初步实验结果显示，该模型在标准的GO功能分类任务中表现良好，在分子功能(MFO)与细胞成分（CCO）两个领域指标超过基线模型。此外，模型展现了较强的泛化能力，验证了其对训练集中未见GO类别的零样本预测能力。

下一阶段，本课题将继续优化模型，进一步提升模型在有监督分类任务与零样本预测任务中的效率与精度。具体而言，准备将多任务训练过程分开进行优化：首先进行对比学习预训练，利用Lora微调技术\cite{devalal2018lora}微调蛋白质编码器与文本编码器，以增强蛋白质与功能描述之间的语义对齐能力；在有监督学习分类任务中，固定经过微调的编码器，仅对分类头进行训练，以进一步提高多标签分类的准确性;在零样本预测任务中，通过计算蛋白质嵌入与功能语义嵌入之间的相似性，实现对未见功能类别的有效预测，从而显著提升模型的泛化性能与实际应用价值。

